% !TeX program = pdflatex
% papers/conversa_vizinhanca.tex — Bandas 1–7 (banal). Teoria escondida.
\documentclass[11pt,a4paper]{article}
\usepackage[utf8]{inputenc}\usepackage[T1]{fontenc}\usepackage[brazil]{babel}
\usepackage[margin=2.4cm]{geometry}
\usepackage[hidelinks]{hyperref}
\newcommand{\code}[1]{\texttt{\small #1}}
% ─────────────────────────────────────────────────────────────────────────────
%  papers/gkcapa.tex — a capa do Reino, mínima, para os papers em `article`.
%
%  O estilo.tex completo é do `report` (títulos de capítulo, skak, …). Aqui fica
%  só o que a capa precisa, para o paper continuar a compilar sozinho com
%  pdflatex. O tradutor próprio (tests/tex) carrega o estilo.tex da raiz / o
%  symlink papers/estilo.tex e avalia o mesmo `\gkcapa`.
% ─────────────────────────────────────────────────────────────────────────────
\definecolor{ouro}{HTML}{B8912F}
\definecolor{ouroclaro}{HTML}{F3E9CE}
\definecolor{tinta}{HTML}{1E1B16}
\definecolor{regua}{HTML}{6E6858}
\providecommand{\gknota}{\fontsize{7.62}{11.05}\selectfont}
\providecommand{\gktexto}{\fontsize{10.50}{15.22}\selectfont}
\providecommand{\gksub}{\fontsize{12.33}{17.87}\selectfont}
\providecommand{\gksec}{\fontsize{14.47}{20.98}\selectfont}
\providecommand{\gkcap}{\fontsize{16.99}{24.63}\selectfont}
\providecommand{\gktit}{\fontsize{23.42}{33.95}\selectfont}
\providecommand{\Prj}{\mathbb{P}}
\providecommand{\Trf}{\mathcal{T}}
\providecommand{\gkcapa}[3]{%
\title{\vspace{-0.6cm}
{\color{ouro}\rule{\textwidth}{1.4pt}}\\[6mm]
\textbf{\gktit\color{tinta} Reino Dourado}\\[3mm]{\gksec\itshape\color{regua} Golden Kingdom}\\[8mm]
{\gkcap\scshape\color{ouro} #1}\\[5mm]
{\gksub\color{regua} #2}\\[2mm]
{\gktexto\color{regua}\itshape #3}\\[7mm]
{\color{ouro}\rule{\textwidth}{0.6pt}}\\[9mm]{\gknota\color{regua}\begin{minipage}{\textwidth}\setlength{\parindent}{0pt}%
\textbf{Aviso de Isenção Algorítmica:} O universo, as leis e as entidades contidas nestas páginas ---
incluindo felinos matriciais, aves de rapina algébricas e amuletos de controle de fluxo --- são
construções de pura ficção formal. Esta obra não descreve a realidade física, não serve como
engenharia reversa do cosmos e não constitui doutrina religiosa. Qualquer semelhança com bugs reais do seu
sistema operacional é mera coincidência (ou falta de reza). Prossiga por sua própria conta e
risco.\end{minipage}}}
\author{}\date{}}

\begin{document}
\gkcapa{Vizinhança Tolerante}
       {bandas 1--7}
       {erro $\to$ abreviação $\to$ oral $\to$ variação $\to$ ruído $\to$ continuidade $\to$ não-colagem}
\maketitle
\begin{abstract}
Corpus de regressão linguística para a histerese operacional:
a assistente responde de forma banal sem ``você quis dizer\ldots''.
A teoria fica no Corpo de Peano; aqui só há situações.
Semente: \code{tools/vizinhanca.sh}. Bateria: \code{tools/bench\_vizinhanca.sh}.
\end{abstract}

\section{nota de escala}
erro; abreviação; variação; continuidade; mudança.
Português coloquial brasileiro. Sem aula de leis.

\section{muzica}
Gosto de música --- deixa o ambiente mais leve.

\section{cafeh}
Café quentinho sempre cai bem.

\section{vc tbm gosta de cafe?}
Gosto sim, especialmente quando está bem quente.

\section{to morrendo de fome}
Então está na hora de arrumar alguma coisa gostosa para comer.

\section{ta frio ai?}
Parece que hoje está pedindo um café quentinho.

\section{cafezim agora ia bem}
Ia mesmo. Vamos a isso.

\section{vc gosta d cafe? kkk}
Gosto sim! Café quente é sempre bem-vindo.

\section{to com fome}
Quer uma sugestão rápida para comer?

\section{sim}
Um sanduíche quente ou uma omelete são opções simples.

\section{sem queijo}
Então dá para fazer uma omelete simples com tomate e ervas.

\section{sem qjo}
Então dá para fazer uma omelete simples com tomate e ervas.

\section{qual é a capital do Chile?}
A capital do Chile é Santiago.

\end{document}
